\section{Grad-CAM và bản đồ đặc trưng}
\label{sec:ch05-grad-cam}

Baseline làm Integrated Gradients phù hợp với attribution trên các tọa độ đầu
vào. Với một ảnh, hàng triệu tọa độ pixel thường quá chi tiết để trả lời câu
hỏi mô hình nhìn vào vùng nào. Grad-CAM chuyển điểm đặt attribution vào lớp tích
chập cuối còn giữ cấu trúc không gian. Mỗi \tn{bản đồ đặc trưng}{feature map}
$A^k$ cho biết một mẫu đã xuất hiện mạnh đến đâu tại từng vị trí, còn các lớp
sau biến tập bản đồ ấy thành điểm số $y^c$ cho lớp mục tiêu $c$.

Chọn lớp tích chập là một đánh đổi. Lớp sâu chứa các mẫu ngữ nghĩa hơn, nhưng
lưới không gian đã thô. Lớp nông giữ chi tiết vị trí hơn, nhưng các bản đồ của
nó chưa gắn chặt với khái niệm mà lớp $c$ biểu thị. Grad-CAM thường dùng lớp
tích chập cuối vì đó là nơi hai yêu cầu này gặp nhau, chứ không vì phép tính
buộc phải dùng đúng lớp ấy~\autocite{p12gradcam}.

Khác với Integrated Gradients, Grad-CAM không so ảnh với một ảnh nền. Nó hỏi
một câu có điều kiện khác: với lớp $c$ đã chọn, bản đồ đặc trưng nào đang nâng
điểm số đó? Câu hỏi phụ thuộc cả đầu ra mục tiêu. Cùng một ảnh có thể cho hai
bản đồ khác nhau khi ta lần lượt giải thích hai lớp, và đó là thuộc tính cần có
của một lời giải thích phân biệt theo lớp.

\index{Grad-CAM|(}%
\index{bản đồ đặc trưng}%
